\section{Open problems}\label{sec:open}

We list what the paper leaves open, in the order of the sections.
\begin{enumerate}[(1)]
\item (The lower half of the matching.) Theorem~\ref{thm:ends}(b) matches the end layer to the
bulk uniformly under (LB$_\plus$). Lemma~\ref{lem:E} gives its upper half under (SRT$_\pm$) for
$\alpha\ne1/2$. Whether the lower half holds under (SRT$_\pm$) alone is open
(Remark~\ref{rem:M0}).
\item (Necessity in the aged law.) Theorem~\ref{thm:local} needs (SRT$_\pm$) for a fresh start.
In the aged decomposition the plus part sees $F^\minus_{N-d}(k-d)$ with the same number
$N-k$ of minus steps for every $d$, so we expect a failure of (SRT$_\minus$) to survive the
average over $d$ and break Theorem~\ref{thm:aging}(a). This is not proved. Whether (SRT$_\pm$)
alone implies (U) when $\alpha\le1/2$, the uniform matching of the aged layer to $L_\theta$,
and the slowly varying factor at $\alpha=1/2$ are open as well.
\item (The line $\alpha_\minus=1/2$.) The uniform matching of Theorem~\ref{thm:dominated}(d) is
stated for $\alpha_\minus\ne1/2$. On the line, both routes of the proof lose a ratio of slowly
varying factors of $T_\minus$, which need not be bounded.
\item (The big-jump zone of the dominated regime.) Between the bulk at the scale $D_N$ and the
lighter end we expect the two-term kernel of~\cite{paperF} with the signs exchanged, and a
window of width $(\log N)^{1-\alpha_\plus}D_N$ when the lighter mean is finite. Nothing is proved
there, and the initial-run theorem of the dominated regime is written only for minus initial
runs (Remark~\ref{rem:dom-initial}).
\item (The exact count.) Proposition~\ref{prop:count} rests on the hypothesis (B), a
first- and second-difference local limit theorem in the bulk. The first-difference half follows
from the proof of Theorem~\ref{thm:local}. The second-difference half needs a second-difference
form of Lemma~\ref{lem:Lstar} on the range of few runs. The count in the dominated regime
is open as well. We expect $M_\plus+M_\minus+m$ modes there, with $m=1$ exactly when the density
$\rho$ has an interior maximum, which for $\mu_\minus<\infty$ is $\alpha_\plus>1/2$, and the
threshold curve in the $(\alpha_\plus,\alpha_\minus)$ square runs from $(\alpha_1,\alpha_1)$ to
$(1/2,1)$. Every exact law we computed agrees with this.
\item (Every $N$.) All counts here are for large $N$. For completely monotone run laws with a
fresh start, every exact law we computed has at most $3$ local maxima in $\{1,\dots,N-1\}$ for
every $N$. For the fair start no such bound holds: there are completely monotone laws, the
same for both signs, whose fair law has $5$ modes at $N=39$ while both fresh laws are
monotone. We know no proof of a bound for every $N$.
\item (The line $\alpha=1$.) When one or both indices equal $1$ the periods may or may not have
a finite mean, and the scales carry slowly varying corrections. The local tools of
Berger~\cite{berger2019ptrf} for the Cauchy domain should give the local law at the scale
$N\ell(N)/m(N)$, with $m$ the truncated mean, and the dominated regime with a lighter law of
index $1$ converges at the rate $1/\log N$. We have not worked out either line.
\end{enumerate}
